\documentclass[../main.tex]{subfiles}
\begin{document}
\subsection{Основные определения}
%==============================================================================

\begin{definition}
В кинематике широко используется идеализированная модель твёрдого тела --- \textit{абсолютно твёрдое тело} (АТТ).
Определяющее свойство этой модели: расстояние между любыми двумя точками тела остаётся неизменным при любом движении:
\[
|AB| = \mathrm{const} \quad \text{для любых точек } A,\, B \text{ тела}.
\]
Это условие \textit{нерастяжимости} является фундаментом всей дальнейшей теории.
Реальные объекты (стальная линейка, маховик, корпус автомобиля) приближённо удовлетворяют ему,
пока упругие деформации пренебрежимо малы. Кусок резины или пружина моделью АТТ не описываются.
\end{definition}

\begin{definition}
Движение АТТ называется \textit{плоскопараллельным}, если в каждый момент времени
векторы скоростей всех точек тела параллельны некоторой неподвижной плоскости~$\Pi$
(плоскости приведения). Следовательно, траектории всех точек лежат в плоскостях,
параллельных~$\Pi$.
\end{definition}

\noindent\textbf{Примеры:} колесо велосипеда, движущегося прямолинейно по дороге; шатун кривошипно-шатунного механизма; шестерня в зубчатой передаче с параллельными осями; палка, скользящая по льду произвольным образом.

\noindent\textbf{Не является плоскопараллельным:} колесо велосипеда при повороте руля; прецессия гироскопа; закрученный мяч, летящий в поле тяжести; винт вертолёта в общем случае полёта.

\begin{definition}
Поскольку все сечения тела, параллельные~$\Pi$, движутся абсолютно одинаково
(их положения в каждый момент совпадают с точностью до переноса вдоль нормали к~$\Pi$),
кинематика всего тела полностью определяется движением одного такого сечения ---
\textit{плоской фигуры} --- в своей плоскости.
Далее мы будем изучать именно движение плоской фигуры, подразумевая, что
результаты автоматически распространяются на всё тело.
\end{definition}

\begin{example}
Рассмотрим тонкий стержень $AB$ длины~$L$, конец которого~$A$ скользит
по горизонтальному полу, а конец~$B$ --- по вертикальной стене
(классическая модель соскальзывающей лестницы). Траектория точки $A$
представляет собой горизонтальную прямую, точки $B$ --- вертикальную
прямую, а любой промежуточной точки $C$ --- эллипс. Поскольку все эти
траектории лежат в параллельных вертикальных плоскостях, движение
стержня является плоскопараллельным. При этом оно \textit{не сводится}
к вращению вокруг неподвижной оси: ни одна точка стержня не покоится
в течение всего движения. Данный пример демонстрирует, что
плоскопараллельное движение --- существенно более общий случай, чем
простое вращение или поступательное перемещение. В дальнейшем мы
покажем, что в каждый момент времени такое движение эквивалентно
мгновенному вращению вокруг \textit{мгновенного центра скоростей},
который сам непрерывно перемещается в пространстве.
\end{example}

%==============================================================================
\subsection{Теорема о перемещении плоской фигуры}
%==============================================================================

\begin{theorem}
Эйлера--Шаля. Любое перемещение плоской фигуры из одного положения в другое
эквивалентно либо повороту вокруг некоторой точки (полюса поворота),
либо параллельному переносу.
\end{theorem}

\begin{proof}
Пусть точки $A$ и $B$ фигуры перешли в $A'$ и $B'$; пусть $A \ne A'$.
Проведём серединные перпендикуляры $l_1$ и $l_2$ к отрезкам $AA'$ и $BB'$.

\noindent\textbf{Случай 1:} $l_1 \cap l_2 = \{O\}$.\\
По построению $|OA| = |OA'|$ и $|OB| = |OB'|$.
Из условия абсолютной твёрдости тела $|AB| = |A'B'|$.
Следовательно, $\triangle OAB \cong \triangle OA'B'$ по трём сторонам.
Так как фигура движется, не вынимаясь из плоскости (как картонный треугольник),
её ориентация сохраняется. Значит, конгруэнтность реализуется
поворотом вокруг $O$, а не зеркальным отражением.

\noindent\textbf{Единственность центра для всей фигуры.}\\
Докажем, что найденная точка $O$ является центром поворота не только
для отрезка $AB$, но и для любой другой точки $P$ фигуры.
Пусть при перемещении $P$ переходит в $P'$. Обозначим через $P^*$
образ точки $P$ при повороте вокруг $O$ на угол $\varphi = \angle AOA'$.
Из сохранения расстояний следует:
$|A'P^*| = |AP| = |A'P'|$ и $|B'P^*| = |BP| = |B'P'|$.
На плоскости существует не более двух точек, одинаково удалённых от $A'$ и $B'$;
они симметричны относительно прямой $A'B'$. Поскольку и исходное движение,
и поворот сохраняют ориентацию фигуры (не «выворачивают» её наизнанку),
точки $P^*$ и $P'$ не могут быть зеркальными отражениями --- они обязаны совпасть.
Следовательно, один и тот же центр $O$ задаёт поворот \textbf{всей} фигуры.

\noindent\textbf{Случай 2:} $l_1 \parallel l_2$, $l_1 \ne l_2$.\\
Тогда отрезки $AA'$ и $BB'$ параллельны. Из условия $|AB| = |A'B'|$
следует, что четырёхугольник $AA'B'B$ --- параллелограмм, откуда
$|AA'| = |BB'|$ и направление перемещения одинаково для всех точек.
Следовательно, перемещение является параллельным переносом.
\end{proof}

\noindent\textit{Важное замечание.}
Полюс $O$ --- это центр \textit{конечного} перемещения фигуры за
некоторый интервал времени. Он, вообще говоря, \textbf{не совпадает}
с мгновенным центром скоростей (МЦС), который будет введён в~\S\,6.
Путать конечный полюс поворота и мгновенный центр скоростей ---
одна из самых распространённых ошибок в кинематике.

\begin{example}
Пусть квадрат переместился так, что его сторона $AB$ перешла в $A'B'$.
По теореме Эйлера--Шаля, центр конечного поворота $O$ находится как
пересечение серединных перпендикуляров к $AA'$ и $BB'$. Даже если
геометрический центр квадрата сместился на 3~см вправо, а сам квадрат
повернулся на $30^\circ$, точка $O$ не совпадает ни с начальным, ни
с конечным положением центра. Её координаты определяются исключительно
геометрическим построением серединных перпендикуляров.
\end{example}

\begin{example}
Треугольник $ABC$ с вершинами $A = (0, 0)$, $B = (2, 0)$, $C = (0, 2)$
совершил перемещение, в результате которого $A' = (2, 0)$, $B' = (2, 2)$,
$C' = (0, 0)$. Найти центр конечного поворота.

Серединный перпендикуляр к $AA'$: середина отрезка --- $(1, 0)$,
$AA'$ горизонтален, значит $l_1$ --- вертикаль $x = 1$.
Серединный перпендикуляр к $BB'$: середина отрезка --- $(2, 1)$,
$BB'$ вертикален, значит $l_2$ --- горизонталь $y = 1$.
Их пересечение даёт полюс $O = (1, 1)$.
Проверка нерастяжимости: $|AC| = 2$, $|A'C'| = \sqrt{(2-0)^2 + (0-0)^2} = 2$. $\checkmark$
Угол поворота: $\angle AOA' = 90^\circ$ (против часовой стрелки). $\checkmark$
\end{example}

%==============================================================================
\subsection{Предельный переход: от конечного перемещения к мгновенному движению}
%==============================================================================

Теорема Эйлера--Шаля описывает \textit{конечные} перемещения:
в какое положение перейдёт фигура за некоторый промежуток времени.
В кинематике же нас интересует \textit{мгновенная} картина:
как тело движется в данный конкретный момент.

Рассмотрим очень малый промежуток времени~$\Delta t$.
За это время тело совершает малое перемещение, которое по теореме Шаля
сводится к повороту на малый угол~$\Delta\varphi$ вокруг некоторого центра~$O(\Delta t)$.
Будем уменьшать $\Delta t$, стремясь к нулю.
При этом из одного геометрического процесса естественным образом возникают две физические величины:
\[
\omega = \lim_{\Delta t \to 0} \frac{\Delta\varphi}{\Delta t} = \dot\varphi
\qquad\text{--- мгновенная угловая скорость;}
\]
\[
C = \lim_{\Delta t \to 0} O(\Delta t)
\qquad\text{--- мгновенный центр скоростей (МЦС).}
\]
Запись $\lim_{\Delta t \to 0}$ означает, что мы исследуем поведение отношения
$\frac{\Delta\varphi}{\Delta t}$ и положения точки $O$ при сколь угодно малых~$\Delta t$.
Если эти значения неограниченно приближаются к определённому числу и точке,
то пределы существуют.

\noindent\textit{Существует ли предел на самом деле?}
Да, при условии, что движение плавное (гладкое): положения точек
и их скорости меняются непрерывно во времени. В этом случае
перпендикуляры к векторам скоростей точек поворачиваются плавно, поэтому их
точка пересечения стремится к фиксированной точке~$C$.
Это и есть мгновенный центр скоростей.

Важно различать два случая параллельных скоростей:
\begin{enumerate}
  \item Если скорости двух точек параллельны и \textit{равны} по модулю и направлению, то перпендикуляры к ним параллельны и не пересекаются в конечной области. Это соответствует $\omega = 0$ (мгновенное поступательное движение), и МЦС уходит на бесконечность.
  \item Если же скорости параллельны, но \textit{различны} по модулю (или направлены в противоположные стороны), то $\omega \ne 0$. В этом случае перпендикуляры совпадают с прямой, соединяющей точки, и МЦС находится на конечном расстоянии (со стороны точки с меньшей скоростью, либо между точками при противоположных направлениях).
\end{enumerate}

Главный вывод: существование МЦС при $\omega \ne 0$
не постулируется отдельно, а \textbf{вытекает из геометрии конечного перемещения}
(теоремы Эйлера--Шаля).

\noindent\textit{Важное замечание.}
МЦС --- понятие строго \textit{мгновенное}.
В следующий момент времени тело будет вращаться вокруг \textit{другой} точки пространства.
Следует различать мгновенный центр скоростей~$C$ и центр конечного поворота~$O$
из теоремы Шаля: это разные точки, определяемые разными способами.

\begin{example}
Представьте, что колесо разгоняется.
Измерим, на какой угол $\Delta\varphi$ оно повернулось за разные промежутки~$\Delta t$:
за $1$~с --- на $6$~рад (средняя скорость $6$~рад/с);
за $0{,}1$~с --- на $0{,}65$~рад (средняя скорость $6{,}5$~рад/с);
за $0{,}01$~с --- на $0{,}066$~рад (средняя скорость $6{,}6$~рад/с).
По мере уменьшения $\Delta t$ среднее значение $\frac{\Delta\varphi}{\Delta t}$
стабилизируется и приближается к некоторому числу. Это число и есть
мгновенная угловая скорость $\omega \approx 6{,}6$~рад/с в данный момент времени.
Точно так же центр малого поворота $O(\Delta t)$ при уменьшении~$\Delta t$
неограниченно приближается к фиксированной точке~$C$ --- мгновенному центру скоростей.
\end{example}

%==============================================================================
\subsection{Угловая скорость}
%==============================================================================

\begin{definition}
Угловой скоростью плоскопараллельного движения называется скорость изменения
угла поворота любой прямой, жёстко связанной с телом:
\[
\omega = \dot\varphi.
\]
Величина $\omega$ едина для всего тела в данный момент времени.
Знак определяется направлением вращения: $\omega > 0$ соответствует
вращению против часовой стрелки, $\omega < 0$ --- по часовой.
Частный случай: поступательное движение, при котором все прямые тела
сохраняют своё направление, соответствует $\omega = 0$.
\end{definition}

\noindent\textbf{Почему $\omega$ одинакова для всех прямых тела.}
Возьмём любые две точки $A$ и $B$ тела. Так как расстояние между ними
неизменно, квадрат длины отрезка постоянен:
\[
|AB|^2 = (x_B - x_A)^2 + (y_B - y_A)^2 = \mathrm{const}.
\]
Продифференцируем равенство по времени:
\[
2(x_B - x_A)(\dot x_B - \dot x_A) + 2(y_B - y_A)(\dot y_B - \dot y_A) = 0.
\]
Векторно это записывается как $\overrightarrow{AB} \cdot (\vec{v}_B - \vec{v}_A) = 0$.
Следовательно, вектор относительной скорости $\vec{v}_B - \vec{v}_A$
всегда перпендикулярен отрезку $AB$. Геометрически это означает,
что точка $B$ мгновенно вращается вокруг $A$.

Отношение модуля относительной скорости к длине отрезка,
\[
\omega = \frac{|\vec{v}_B - \vec{v}_A|}{|AB|},
\]
не зависит от выбора пары точек $A$ и $B$ (это прямое следствие жёсткой
связи всех точек тела). Поэтому $\omega$ является характеристикой движения
всего тела, а не отдельного отрезка.

\noindent\textbf{Угловая скорость и система отсчёта.}
Возникает естественный вопрос: изменится ли $\omega$, если перейти к другой системе отсчёта? Рассмотрим простейший, но фундаментальный случай. Пусть система отсчёта $S'$ движется относительно лабораторной системы $S$ поступательно со скоростью $\vec{u}$ (без вращения). Скорость любой точки $P$ тела в системе $S'$ связана со скоростью в $S$ правилом сложения: $\vec{v}'_P = \vec{v}_P - \vec{u}$.

Вычислим разность скоростей двух произвольных точек $A$ и $B$ жёсткого тела в новой системе:
\[
\vec{v}'_B - \vec{v}'_A = (\vec{v}_B - \vec{u}) - (\vec{v}_A - \vec{u}) = \vec{v}_B - \vec{v}_A.
\]
Скорость переноса $\vec{u}$ полностью сокращается. Следовательно, угловая скорость в системе $S'$ вычисляется точно так же:
\[
\omega' = \frac{|\vec{v}'_B - \vec{v}'_A|}{|AB|} = \frac{|\vec{v}_B - \vec{v}_A|}{|AB|} = \omega.
\]
Мы получили важный результат: угловая скорость плоскопараллельного движения \textbf{инвариантна} относительно перехода к любой поступательно движущейся системе отсчёта. Она характеризует вращение тела само по себе и не зависит от движения наблюдателя, если только его система не вращается. Этот факт строго обосновывает уже знакомое свойство: выбор любой точки $A$ в качестве «базы» в формуле поля скоростей эквивалентен переходу в поступательно движущуюся систему, поэтому $\omega$ остаётся неизменной.

\begin{example}
Стержень $AB$ длины $L = 1$~м вращается вокруг неподвижного конца~$A$
с угловой скоростью $\omega = 2$~рад/с. Найдём скорости точки~$B$
и точки~$C$, удалённой от~$A$ на $0{,}4$~м.

Поскольку точка $A$ покоится, скорость любой точки стержня определяется
формулой $v = \omega r$, где $r$ --- расстояние до оси вращения,
а направление скорости перпендикулярно стержню. Подставляем числа:
\[
v_B = 2 \cdot 1 = 2~\text{м/с}, \qquad
v_C = 2 \cdot 0{,}4 = 0{,}8~\text{м/с}.
\]
\end{example}

\begin{example}
Стержень $AB$ длины $L$ движется, скользя концами: $A$ --- по горизонтальному
полу, $B$ --- по вертикальной стене. В некоторый момент стержень образует
с полом угол $\alpha = 30^\circ$, а точка $A$ движется вправо со скоростью
$u = 2$~м/с. Найдём $\omega$ и скорость $v_B$.

Направим оси: $x$ вправо, $y$ вверх. Тогда $\vec{v}_A = (u, 0)$,
$\vec{v}_B = (0, -v_B)$ (точка $B$ скользит вниз).
Вектор $\overrightarrow{AB}$ направлен от $A$ к $B$, его единичный вектор:
$\vec{e}_{AB} = (-\cos\alpha, \sin\alpha)$.

Условие нерастяжимости требует, чтобы разность скоростей
$\vec{v}_B - \vec{v}_A = (-u, -v_B)$ была перпендикулярна $\overrightarrow{AB}$.
Их скалярное произведение равно нулю:
\[
(-u)(-\cos\alpha) + (-v_B)(\sin\alpha) = 0
\quad\Longrightarrow\quad
v_B = u \cot\alpha.
\]
При $\alpha = 30^\circ$ получаем $v_B = 2\sqrt{3} \approx 3{,}5$~м/с.
Модуль относительной скорости:
$|\vec{v}_B - \vec{v}_A| = \sqrt{u^2 + v_B^2} = 4$~м/с.
Угловая скорость стержня:
\[
\omega = -\frac{4}{L}~\text{рад/с}.
\]
Отрицательный знак указывает на вращение по часовой стрелке.
\end{example}

%==============================================================================
\subsection{Поле скоростей плоской фигуры}
%==============================================================================

Из предыдущего раздела следует, что для любых двух точек $A$ и $B$ абсолютно твёрдого тела
относительная скорость $\vec{v}_B - \vec{v}_A$ перпендикулярна отрезку $AB$,
а её модуль равен $\omega |AB|$. Это означает, что в данный момент точка $B$
движется так, как если бы она мгновенно вращалась вокруг $A$ с угловой скоростью $\omega$.

Введём декартову систему координат. Пусть $\overrightarrow{AB} = (x_{AB}, y_{AB})$,
где $x_{AB} = x_B - x_A$, $y_{AB} = y_B - y_A$.
Вектор, перпендикулярный $\overrightarrow{AB}$ и повёрнутый на $90^\circ$
против часовой стрелки, имеет координаты $(-y_{AB}, x_{AB})$.
Поэтому относительная скорость записывается как
$\vec{v}_B - \vec{v}_A = \omega (-y_{AB}, x_{AB})$.
Отсюда получаем основную формулу поля скоростей:
\[
\vec{v}_B = \vec{v}_A + \omega (-y_{AB},\; x_{AB}),
\]
или в покомпонентной записи:
\[
v_{Bx} = v_{Ax} - \omega y_{AB}, \qquad
v_{By} = v_{Ay} + \omega x_{AB}.
\]
Тем самым всё поле скоростей фигуры в данный момент полностью определяется
скоростью всего одной точки $\vec{v}_A$ и единой угловой скоростью $\omega$.
(В трёхмерной записи это соответствует формуле Эйлера
$\vec{v}_B = \vec{v}_A + \vec{\omega} \times \vec{r}_{AB}$,
где вектор $\vec{\omega}$ направлен по нормали к плоскости движения.)

\noindent\textbf{Приём~1. Теорема о проекциях скоростей.}
Условие перпендикулярности $(\vec{v}_B - \vec{v}_A) \perp \overrightarrow{AB}$
можно переформулировать через скалярное произведение:
\[
(\vec{v}_B - \vec{v}_A) \cdot \overrightarrow{AB} = 0
\quad\Longleftrightarrow\quad
\text{пр}_{AB}(\vec{v}_B) = \text{пр}_{AB}(\vec{v}_A).
\]
Это означает, что \textit{проекции скоростей любых двух точек тела
на прямую, их соединяющую, равны}.
Данный факт является прямым следствием нерастяжимости тела
и не требует знания ни $\omega$, ни координат точек.
На практике теорема о проекциях часто оказывается быстрее
и надёжнее, чем полная векторная формула, особенно при работе
со стержнями, нитями и ползунами.

\begin{example}
Нить горизонтально тянут вправо со скоростью $v$, разматывая её
с катушки, которая катится по плоскости без скольжения.
Внутренний радиус катушки $r$, внешний --- $R$.
Нить сматывается с верхней части внутреннего цилиндра.
Найдём скорость центра катушки $u$.

При качении без скольжения точка касания с опорой покоится,
поэтому угловая скорость связана со скоростью центра соотношением
$\omega = \frac{u}{R}$. Точка $K$, в которой нить отрывается от катушки, находится
строго над центром $O$ на расстоянии $r$.
Применяя формулу поля скоростей относительно точки $O$, получаем:
скорость $K$ складывается из переносной скорости $u$ и вращательной $\omega r$,
причём при качении вправо оба вектора сонаправлены. Следовательно,
$v = u + \omega r = u + \frac{u}{R}r = u \frac{R+r}{R}$,
откуда
\[
u = v \frac{R}{R+r}.
\]
\end{example}

\begin{example}
Кривошип $OA$ длины $a = 0{,}2$~м вращается вокруг неподвижной точки $O$
с угловой скоростью $\omega_0 = 5$~рад/с.
Шатун $AB$ длины $b = 0{,}5$~м соединяет точку $A$ с ползуном $B$,
который может двигаться только горизонтально.
В рассматриваемый момент кривошип $OA$ горизонтален и направлен вправо,
а шатун $AB$ также горизонтален. Найдём скорость ползуна $v_B$
и угловую скорость шатуна $\omega_{AB}$.

Скорость точки $A$ равна $v_A = \omega_0 a = 1$~м/с и направлена
вертикально вверх (перпендикулярно горизонтальному кривошипу).
Ползун $B$ может двигаться только вдоль горизонтали.
Применим теорему о проекциях скоростей на прямую $AB$ (горизонтальную).
Проекция $\vec{v}_A$ на горизонталь равна нулю, значит,
и проекция $\vec{v}_B$ должна быть нулевой. Но $\vec{v}_B$ сам направлен
горизонтально, следовательно, $v_B = 0$.
В данный момент ползун мгновенно покоится (это положение называется
\textit{мёртвой точкой} механизма).

Для нахождения $\omega_{AB}$ воспользуемся формулой поля скоростей.
Направим ось $y$ вертикально вверх. Тогда $v_{Ay} = 1$~м/с, $v_{By} = 0$,
а координата $x_{AB} = b = 0{,}5$~м.
Из соотношения $v_{By} = v_{Ay} + \omega_{AB} x_{AB}$ получаем:
\[
0 = 1 + \omega_{AB} \cdot 0{,}5 \quad\Longrightarrow\quad
\omega_{AB} = -2~\text{рад/с}.
\]
Отрицательный знак указывает на вращение шатуна по часовой стрелке.
\end{example}

%==============================================================================
\subsection{Мгновенный центр скоростей}
%==============================================================================

\begin{definition}
Мгновенным центром скоростей (МЦС) называется точка плоской фигуры
(или её геометрического продолжения), скорость которой в данный момент
времени равна нулю. МЦС существует и единственен при условии $\omega \ne 0$.
Это строго мгновенное понятие: в следующий момент времени МЦС, как правило,
переместится в другую точку пространства.
\end{definition}

\noindent\textbf{Как найти МЦС.}
Из формулы поля скоростей следует, что скорость любой точки $A$ направлена
перпендикулярно отрезку $CA$, соединяющему её с МЦС $C$. Следовательно,
если известно направление скорости точки $A$, то МЦС обязательно лежит на
прямой, проведённой через $A$ перпендикулярно $\vec{v}_A$. Чтобы однозначно
найти точку $C$, достаточно построить такие перпендикуляры к скоростям
двух любых точек $A$ и $B$ (при условии, что скорости не параллельны).
Их пересечение и есть искомый МЦС.

\noindent\textbf{Поле скоростей через МЦС.}
Если принять МЦС за полюс, формула поля скоростей приобретает вид чистого
вращения вокруг неподвижной оси:
\[
v_P = \omega \cdot |CP|, \qquad \vec{v}_P \perp CP,
\]
причём вектор скорости направлен перпендикулярно радиусу $CP$ в сторону
вращения фигуры. Это позволяет вычислять скорости точек без векторного
сложения: достаточно знать расстояние до МЦС, общую угловую скорость
и геометрическое направление.

\begin{example}
Колесо радиуса $R$ катится без скольжения вправо; скорость центра $O$
равна $v$. По условию качения точка касания $C$ с дорогой мгновенно
покоится, значит, она и есть МЦС. Угловая скорость колеса
$\omega = \frac{v_O}{R} = \frac{v}{R}$. Теперь найдём скорости других точек,
рассматривая колесо как мгновенно вращающееся вокруг $C$:
верхняя точка $A$ находится на расстоянии $2R$ от $C$, поэтому
$v_A = \omega \cdot 2R = 2v$ (направлена горизонтально вправо);
правая точка $B$ на конце горизонтального диаметра удалена от $C$
на $R\sqrt{2}$, значит, $v_B = v\sqrt{2}$, а её скорость направлена
под углом $45^\circ$ вверх-вправо (перпендикулярно отрезку $CB$).
Аналогично, левая точка $D$ имеет скорость $v_D = v\sqrt{2}$,
направленную под $45^\circ$ вверх-влево.
\end{example}

\begin{example}
Стержень $AB$ длины $L$ скользит концами: $A$ по полу, $B$ по стене.
В момент, когда стержень образует с полом угол $\alpha = 30^\circ$,
точка $A$ движется вправо со скоростью $u$. Найдём МЦС, $\omega$
и скорость середины $C$ стержня.

Скорость $\vec{v}_A$ горизонтальна, поэтому перпендикуляр к ней через $A$ ---
вертикаль. Скорость $\vec{v}_B$ вертикальна, перпендикуляр к ней через $B$ ---
горизонталь. Их пересечение даёт точку $M$ (МЦС). Из геометрии
прямоугольного треугольника $MAB$ катет $|MA| = L\sin\alpha$,
катет $|MB| = L\cos\alpha$. Угловая скорость стержня
$\omega = \frac{u}{|MA|} = \frac{u}{L\sin\alpha}$.
Середина стержня $C$ равноудалена от вершин прямоугольного треугольника,
значит, $|MC| = \frac{L}{2}$. Скорость середины
$v_C = \omega \cdot \frac{L}{2} = \frac{u}{2\sin\alpha}$.
При $u = 2$ м/с и $\alpha = 30^\circ$ получаем $\omega = \frac{4}{L}$ рад/с,
$v_C = 2$ м/с. Направление $\vec{v}_C$ перпендикулярно отрезку $MC$.
\end{example}

\begin{example}
Ролик радиуса $R$ зажат между двумя горизонтальными досками.
Верхняя доска движется вправо со скоростью $v_1 = 4$ м/с,
нижняя --- вправо со скоростью $v_2 = 2$ м/с. Найдём $\omega$
и скорость оси ролика $O$.

Точки касания ролика с досками имеют те же скорости, что и доски.
Обе скорости горизонтальны и сонаправлены, поэтому перпендикуляры
к ним вертикальны и параллельны. В таком случае МЦС лежит на прямой,
соединяющей точки касания (вертикальный диаметр ролика).
Из линейного распределения скоростей по вертикали или из подобия
треугольников следует, что ось $O$ находится посередине между
точками с разными скоростями, поэтому её скорость --- среднее
арифметическое: $v_O = \frac{v_1 + v_2}{2} = 3$ м/с.
Угловая скорость определяется разностью скоростей на диаметре:
$\omega R = \frac{v_1 - v_2}{2} \Rightarrow \omega = \frac{v_1 - v_2}{2R} = \frac{1}{R}$ рад/с
(вращение по часовой стрелке).
\end{example}

\noindent\textbf{Типовые случаи нахождения МЦС.}
На практике часто встречаются стандартные конфигурации, которые полезно
запомнить. При качении без скольжения по неподвижной поверхности МЦС
совпадает с точкой касания. Если известны направления скоростей двух
точек, МЦС находится пересечением перпендикуляров к этим векторам.
Если скорости двух точек параллельны, но различны по модулю, МЦС лежит
на прямой, соединяющей точки, за пределами отрезка --- со стороны точки
с меньшей скоростью. Его положение находится из подобия:
$\frac{|CA|}{|CB|} = \frac{v_A}{v_B}$. Если же скорости параллельны и направлены
в противоположные стороны, МЦС лежит \textit{между} точками.
Если скорости параллельны и равны по модулю, движение мгновенно
поступательное, $\omega = 0$, и МЦС уходит в бесконечность.

\noindent\textbf{Центроиды.}
Поскольку МЦС перемещается со временем, он описывает две кривые:
траекторию в пространстве (\textit{неподвижную центроиду}) и траекторию
относительно самого тела (\textit{подвижную центроиду}). Геометрически
плоскопараллельное движение эквивалентно качению подвижной центроиды
по неподвижной без проскальзывания. Для колеса, катящегося по прямой,
неподвижная центроида --- это сама дорога, а подвижная --- обод колеса.

\noindent\textit{Важное замечание об ускорениях.}
Метод МЦС применим \textit{только к скоростям}. Для ускорений существует
свой мгновенный центр, но его положение не совпадает с МЦС, а в формулах
появляется дополнительное центростремительное слагаемое $\omega^2 r$.
Попытка найти ускорения по прямой аналогии с $v = \omega r$ ---
распространённая и серьёзная ошибка.

%==============================================================================
\subsection{Ускорения точек твёрдого тела}
%==============================================================================

Для нахождения ускорений продифференцируем формулу поля скоростей по времени. Запишем её в покомпонентном виде:
\[
v_{Bx} = v_{Ax} - \omega\, y_{AB}, \qquad
v_{By} = v_{Ay} + \omega\, x_{AB},
\]
где $x_{AB} = x_B - x_A$ и $y_{AB} = y_B - y_A$ сами зависят от времени, так как точки тела движутся. Применяя правило дифференцирования произведения и вводя угловое ускорение $\varepsilon = \dot\omega$, получаем для первой координаты:
\[
a_{Bx} = a_{Ax} - \varepsilon\, y_{AB} - \omega\,\dot y_{AB}.
\]
Производную $\dot y_{AB}$ легко выразить через скорости: $\dot y_{AB} = \dot y_B - \dot y_A = v_{By} - v_{Ay}$. Но из той же формулы поля скоростей мы знаем, что $v_{By} - v_{Ay} = \omega\, x_{AB}$. Подставляем это выражение вместо $\dot y_{AB}$:
\[
a_{Bx} = a_{Ax} - \varepsilon\, y_{AB} - \omega^2\, x_{AB}.
\]
Аналогично дифференцируем вторую координату, учитывая, что $\dot x_{AB} = v_{Bx} - v_{Ax} = -\omega\, y_{AB}$:
\[
a_{By} = a_{Ay} + \varepsilon\, x_{AB} - \omega^2\, y_{AB}.
\]
Объединяя компоненты, приходим к компактной векторной записи формулы Эйлера для ускорений:
\[
\vec{a}_B = \vec{a}_A
+ \varepsilon\,(-y_{AB},\; x_{AB})
- \omega^2\,(x_{AB},\; y_{AB}).
\]

Полученное выражение раскладывает ускорение точки $B$ на три геометрически различимых слагаемых. Первое, $\vec{a}_A$, --- ускорение полюса, выбранного нами в качестве базы отсчёта. Второе слагаемое, $\varepsilon(-y_{AB},\, x_{AB})$, представляет собой вектор, перпендикулярный отрезку $AB$ и повёрнутый на $90^\circ$ против часовой стрелки. Его модуль равен $\varepsilon|AB|$. Это \textit{касательное (тангенциальное) ускорение}, которое возникает исключительно за счёт изменения угловой скорости. Если вращение равномерное ($\varepsilon = 0$), данный член обращается в нуль. Третье слагаемое, $-\omega^2(x_{AB},\, y_{AB})$, направлено строго от точки $B$ к точке $A$, а его модуль равен $\omega^2|AB|$. Это \textit{центростремительное (нормальное) ускорение}. Оно присутствует всегда, когда тело вращается ($\omega \ne 0$), даже если угловая скорость постоянна, и отвечает за искривление траектории точки относительно полюса.

\begin{example}
\textbf{Ускорение точки касания катящегося диска.} Этот пример демонстрирует принципиальное отличие ускорений от скоростей. Диск радиуса $R$ катится без скольжения вправо. Скорость центра $O$ равна $v$, а его ускорение $\vec{a}_O$ направлено горизонтально и равно $a_O$. При качении без скольжения мгновенная угловая скорость и угловое ускорение связаны со скоростью и ускорением центра соотношениями $|\omega| = v/R$ и $|\varepsilon| = a_O/R$.

Точка касания $C$ с опорой является мгновенным центром скоростей, поэтому $\vec{v}_C = \vec{0}$. Найдём ускорение этой точки, выбрав в качестве полюса центр диска: $A = O$, $B = C$. Вектор $\overrightarrow{OC}$ направлен вертикально вниз, следовательно, в нашей системе координат ($x$ вправо, $y$ вверх) его компоненты: $x_{OC} = 0$, $y_{OC} = -R$.

Важно аккуратно учесть знаки. При движении вправо вращение диска происходит по часовой стрелке. Если принять направление против часовой стрелки за положительное, то алгебраические значения угловой скорости и углового ускорения будут отрицательными: $\omega = -v/R$, $\varepsilon = -a_O/R$. Подставляем в формулу Эйлера:
\[
a_{Cx} = a_{Ox} - \varepsilon\, y_{OC} - \omega^2\, x_{OC}
= a_O - \left(-\frac{a_O}{R}\right)(-R) - \omega^2\cdot 0
= a_O - a_O = 0,
\]
\[
a_{Cy} = a_{Oy} + \varepsilon\, x_{OC} - \omega^2\, y_{OC}
= 0 + \varepsilon\cdot 0 - \omega^2(-R)
= \omega^2 R = \frac{v^2}{R}.
\]
Итак, $\vec{a}_C = (0,\; v^2/R)$. Ускорение точки касания направлено строго \textit{вертикально вверх} (к центру диска) и по величине равно $v^2/R$.

Этот результат имеет фундаментальное значение: точка касания в данный момент покоится ($\vec{v}_C = \vec{0}$), но её ускорение \textbf{не равно нулю}. Нулевая скорость мгновенного характера не гарантирует нулевого ускорения. Именно поэтому методы, работавшие для скоростей (например, вращение вокруг МЦС), нельзя механически переносить на ускорения.
\end{example}

Продифференцируем по времени тождество $(\vec{v}_B - \vec{v}_A)\cdot\overrightarrow{AB} = 0$, которое выражает условие нерастяжимости тела. Применяя правило дифференцирования скалярного произведения, получаем:
\[
(\vec{a}_B - \vec{a}_A)\cdot\overrightarrow{AB}
+
(\vec{v}_B - \vec{v}_A)\cdot\frac{d}{dt}\overrightarrow{AB}
= 0.
\]
Производная радиус-вектора $\frac{d}{dt}\overrightarrow{AB}$ равна разности скоростей концов: $\vec{v}_B - \vec{v}_A$. Следовательно, второе слагаемое равно квадрату модуля относительной скорости: $(\vec{v}_B - \vec{v}_A)^2 = |\vec{v}_B - \vec{v}_A|^2 = \omega^2|AB|^2$. Подставляем и делим обе части на $|AB|$:
\[
(\vec{a}_B - \vec{a}_A)\cdot\frac{\overrightarrow{AB}}{|AB|} = -\omega^2|AB|.
\]
Скалярное произведение $(\vec{a}_B - \vec{a}_A)\cdot\frac{\overrightarrow{AB}}{|AB|}$ --- это разность проекций ускорений на прямую $AB$. Обозначая проекции через $a_\parallel$, записываем окончательный результат:
\[
a_{B\parallel} - a_{A\parallel} = -\omega^2|AB|.
\]
Для скоростей разность проекций на соединяющую прямую равнялась нулю. Для ускорений появляется член $-\omega^2|AB|$, который физически отражает центростремительное «стягивание» точек тела относительно друг друга. Эта теорема позволяет быстро находить проекции ускорений без полного разложения на касательную и нормальную составляющие, что особенно удобно при анализе сложных механизмов.

Точку тела (или его продолжения), в которой ускорение в данный момент равно нулю, называют \textit{мгновенным центром ускорений} (МЦУ). Из формулы Эйлера следует, что при условии $\omega^2 + \varepsilon^2 \ne 0$ система уравнений $\vec{a}_P = \vec{0}$ всегда имеет единственное решение, то есть МЦУ существует. Однако поле ускорений вокруг этой точки устроено существенно сложнее, чем поле скоростей вокруг МЦС: ускорения точек не направлены перпендикулярно радиусу и не пропорциональны расстоянию до МЦУ из-за сложного наложения касательного и центростремительного слагаемых. В результате геометрическая картина теряет наглядность, и на практике МЦУ почти не применяется: для нахождения ускорений значительно эффективнее использовать общую формулу Эйлера или теорему о проекциях.

\end{document}